% GENERATED FILE. Edit canonical lessons, then run npm run generate:books.
% canonical-source-hash: fnv1a64:3a2ee3410fdc7757
% canonical-lessons: LA-C527-requiesco, LA-C527-requies, LA-C527-dormito, LA-C527-insomnia, LA-C527-recumbo

\chapter{Rest and sleep}
\label{ch:la-rest-and-sleep}
\hlchaptermodality{527}

\begin{chapteropening}
\textbf{By the end of this chapter:} \emph{I can talk about rest, dozing and sleeplessness.}

Complete the last lesson of chapter 527: I can talk about rest, dozing and sleeplessness.
\end{chapteropening}

\section[requiēscō, requiēscere]{requiēscō, requiēscere — to rest, take a rest}

\label{lesson:LA-C527-requiesco}



\begin{quote}
Before the new one: say the Latin for to refresh, then the Latin for a draught.
\end{quote}

\subsection*{You'll want to know: requiēscō, requiēscere}
\textbf{requiēscō, requiēscere} — \textquotedblleft{}to rest, take a rest\textquotedblright{}.

\begin{quote}
Post labōrem in hortō requiēscimus.
\end{quote}

\begin{quote}
\emph{After work we rest in the garden.}
\end{quote}

\emph{requiēscō, requiēscere, requiēvī, requiētum}. On tombs: \emph{requiēscat in pāce} — may he rest in peace.

\begin{cousinweb}
\textbf{requiēscō} is built on \emph{re-} and \emph{quiēscō} \textquotedblleft{}to rest\textquotedblright{} (Lewis \& Short), from \emph{quiēs} \textquotedblleft{}quiet\textquotedblright{}.
\end{cousinweb}

\begin{grammarlens}[title={What you've built}]
One more word for telling, explaining and describing at B1.
\end{grammarlens}

\subsection*{Your turn}
\begin{itemize}
\raggedright
  \item \emph{Say it:} \emph{requiēscō, requiēscere}
  \item \emph{Say it:} \emph{requiēscō, requiēscere}, once more
  \item \emph{Say it:} a sentence of your own with \emph{requiēscō, requiēscere}
  \item \emph{Recall:} say the Latin for to refresh, then the Latin for a draught, then say \emph{requiēscō, requiēscere} again
\end{itemize}

\begin{tcolorbox}[breakable,colback=teal!4,colframe=teal!35!black,title={Before you move on}]
What does requiēscō, requiēscere mean? (\textquotedblleft{}To rest, take a rest\textquotedblright{}.) Say it once more.
\end{tcolorbox}
\section[requiēs, requiētis]{requiēs, requiētis — rest, repose}

\label{lesson:LA-C527-requies}



\begin{quote}
Before the new one: say the Latin for a draught, then the Latin for to rest.
\end{quote}

\subsection*{You'll want to know: requiēs, requiētis}
\textbf{requiēs, requiētis} — \textquotedblleft{}rest, repose\textquotedblright{}.

\begin{quote}
Nautae nūllam requiem habent.
\end{quote}

\begin{quote}
\emph{The sailors have no rest.}
\end{quote}

Feminine, third declension; accusative \emph{requiem}. English \emph{requiem}, from the first word of a prayer for the dead.

\begin{cousinweb}
\textbf{requiēs} is built on \emph{re-} and \emph{quiēs} \textquotedblleft{}rest, quiet\textquotedblright{} (EtymDB; Lewis \& Short), which EtymDB traces to a Proto-Indo-European word for \textquotedblleft{}rest\textquotedblright{}. English \emph{requiem} comes from its accusative.
\end{cousinweb}

\begin{grammarlens}[title={What you've built}]
One more word for telling, explaining and describing at B1.
\end{grammarlens}

\subsection*{Your turn}
\begin{itemize}
\raggedright
  \item \emph{Say it:} \emph{requiēs, requiētis}
  \item \emph{Say it:} \emph{requiēs, requiētis}, once more
  \item \emph{Say it:} a sentence of your own with \emph{requiēs, requiētis}
  \item \emph{Recall:} say the Latin for a draught, then the Latin for to rest, then say \emph{requiēs, requiētis} again
\end{itemize}

\begin{tcolorbox}[breakable,colback=teal!4,colframe=teal!35!black,title={Before you move on}]
What does requiēs, requiētis mean? (\textquotedblleft{}Rest, repose\textquotedblright{}.) Say it once more.
\end{tcolorbox}
\section[dormītō, dormītāre]{dormītō, dormītāre — to doze, be sleepy}

\label{lesson:LA-C527-dormito}



\begin{quote}
Before the new one: say the Latin for to rest, then the Latin for rest.
\end{quote}

\subsection*{You'll want to know: dormītō, dormītāre}
\textbf{dormītō, dormītāre} — \textquotedblleft{}to doze, be sleepy\textquotedblright{}.

\begin{quote}
Avus post prandium in sellā dormītat.
\end{quote}

\begin{quote}
\emph{Grandfather dozes in his chair after lunch.}
\end{quote}

\emph{dormītō, dormītāre, dormītāvī}: the frequentative of \emph{dormiō}.

\begin{cousinweb}
\textbf{dormītō} is the frequentative of \emph{dormiō} \textquotedblleft{}to sleep\textquotedblright{} (EtymDB; Lewis \& Short). English \emph{dormitory} comes from \emph{dormītōrium}, built on the same verb.
\end{cousinweb}

\begin{grammarlens}[title={What you've built}]
One more word for telling, explaining and describing at B1.
\end{grammarlens}

\subsection*{Your turn}
\begin{itemize}
\raggedright
  \item \emph{Say it:} \emph{dormītō, dormītāre}
  \item \emph{Say it:} \emph{dormītō, dormītāre}, once more
  \item \emph{Say it:} a sentence of your own with \emph{dormītō, dormītāre}
  \item \emph{Recall:} say the Latin for to rest, then the Latin for rest, then say \emph{dormītō, dormītāre} again
\end{itemize}

\begin{tcolorbox}[breakable,colback=teal!4,colframe=teal!35!black,title={Before you move on}]
What does dormītō, dormītāre mean? (\textquotedblleft{}To doze, be sleepy\textquotedblright{}.) Say it once more.
\end{tcolorbox}
\section[īnsomnia, īnsomniae]{īnsomnia, īnsomniae — sleeplessness}

\label{lesson:LA-C527-insomnia}



\begin{quote}
Before the new one: say the Latin for rest, then the Latin for to doze.
\end{quote}

\subsection*{You'll want to know: īnsomnia, īnsomniae}
\textbf{īnsomnia, īnsomniae} — \textquotedblleft{}sleeplessness\textquotedblright{}.

\begin{quote}
Propter cūrās īnsomniā labōrō.
\end{quote}

\begin{quote}
\emph{Because of my worries I suffer from sleeplessness.}
\end{quote}

Feminine, first declension, from \emph{īnsomnis} — sleepless. English \emph{insomnia}.

\begin{cousinweb}
\textbf{īnsomnia} is built on \emph{īnsomnis} \textquotedblleft{}sleepless\textquotedblright{}, from \emph{in-} \textquotedblleft{}not\textquotedblright{} and \emph{somnus} \textquotedblleft{}sleep\textquotedblright{} (Lewis \& Short; EtymDB). English \emph{insomnia} comes from it.
\end{cousinweb}

\begin{grammarlens}[title={What you've built}]
One more word for telling, explaining and describing at B1.
\end{grammarlens}

\subsection*{Your turn}
\begin{itemize}
\raggedright
  \item \emph{Say it:} \emph{īnsomnia, īnsomniae}
  \item \emph{Say it:} \emph{īnsomnia, īnsomniae}, once more
  \item \emph{Say it:} a sentence of your own with \emph{īnsomnia, īnsomniae}
  \item \emph{Recall:} say the Latin for rest, then the Latin for to doze, then say \emph{īnsomnia, īnsomniae} again
\end{itemize}

\begin{tcolorbox}[breakable,colback=teal!4,colframe=teal!35!black,title={Before you move on}]
What does īnsomnia, īnsomniae mean? (\textquotedblleft{}Sleeplessness\textquotedblright{}.) Say it once more.
\end{tcolorbox}
\section[recumbō, recumbere]{recumbō, recumbere — to recline, lie down}

\label{lesson:LA-C527-recumbo}



\begin{quote}
Before the new one: say the Latin for to doze, then the Latin for sleeplessness.
\end{quote}

\subsection*{You'll want to know: recumbō, recumbere}
\textbf{recumbō, recumbere} — \textquotedblleft{}to recline, lie down\textquotedblright{}.

\begin{quote}
Hospitēs ad cēnam recumbunt.
\end{quote}

\begin{quote}
\emph{The guests recline at dinner.}
\end{quote}

\emph{recumbō, recumbere, recubuī}. Romans lay on couches to eat. English \emph{recumbent}.

\begin{cousinweb}
\textbf{recumbō} is built on \emph{re-} \textquotedblleft{}back\textquotedblright{} and \emph{-cumbō}, a form of \emph{cubō} \textquotedblleft{}to lie down\textquotedblright{} (Lewis \& Short). English \emph{recumbent} comes from it.
\end{cousinweb}

\begin{grammarlens}[title={What you've built}]
One more word for telling, explaining and describing at B1.
\end{grammarlens}

\subsection*{Your turn}
\begin{itemize}
\raggedright
  \item \emph{Say it:} \emph{recumbō, recumbere}
  \item \emph{Say it:} \emph{recumbō, recumbere}, once more
  \item \emph{Say it:} a sentence of your own with \emph{recumbō, recumbere}
  \item \emph{Recall:} say the Latin for to doze, then the Latin for sleeplessness, then say \emph{recumbō, recumbere} again
\end{itemize}

\begin{tcolorbox}[breakable,colback=teal!4,colframe=teal!35!black,title={Before you move on}]
What does recumbō, recumbere mean? (\textquotedblleft{}To recline, lie down\textquotedblright{}.) Say it once more.
\end{tcolorbox}
